% POL-2000 · Travail de mi-session (A26) · consignes aux étudiant.e.s
% H24 instructions by Adrien Cloutier, kept nearly verbatim. Build: ./build.sh (xelatex, twice)
\documentclass[11pt]{article}
\usepackage{fontspec}
\usepackage[letterpaper,top=0.6in,bottom=0.75in,left=0.8in,right=0.8in,footskip=0.4in]{geometry}
\usepackage{fancyhdr}
\usepackage{lastpage}
\usepackage{xcolor}
\usepackage{graphicx}
\usepackage{tabularx}
\usepackage{array}
\usepackage{enumitem}
\usepackage[most]{tcolorbox}
\usepackage[hidelinks]{hyperref}
\usepackage{parskip}
\usepackage{newunicodechar}
\usepackage{needspace}
\usepackage{ragged2e}
% Espace fine insécable : XeTeX coupait la ligne dessus.
\newunicodechar{ }{\nobreak\hspace{0.16em}\nobreak}

\definecolor{rouge}{HTML}{C6102E}
\definecolor{encre}{HTML}{121314}
\definecolor{papier}{HTML}{F4F4F0}
\definecolor{gris}{HTML}{6B6B6B}

\setmainfont{IBMPlexSans}[Extension=.otf,UprightFont=*-Regular,BoldFont=*-SemiBold,ItalicFont=*-Italic]
\setmonofont{IBMPlexMono}[Extension=.otf,UprightFont=*-Regular,BoldFont=*-SemiBold,ItalicFont=*-Italic]
\newfontfamily\titre{IBMPlexMono}[Extension=.otf,UprightFont=*-SemiBold,BoldFont=*-Bold]

\hypersetup{pdftitle={POL-2000 · Travail de mi-session · Consignes},pdfauthor={Laurence-Olivier M. Foisy}}
\color{encre}
\frenchspacing
\setlength{\parskip}{4pt}
\linespread{1.04}

\pagestyle{fancy}
\fancyhf{}
\renewcommand{\headrulewidth}{0pt}
\fancyfoot[L]{\small\ttfamily\color{gris}POL-2000 · Travail de mi-session · Automne 2026}
\fancyfoot[R]{\small\ttfamily\color{gris}\thepage{} / \pageref*{LastPage}}

% Titre de section : mono, filet de 2 pt
\newcommand{\etape}[2]{%
  \par\vspace{10pt}\Needspace{4\baselineskip}%
  {\titre\small\color{rouge}#1}\par\vspace{-2pt}%
  {\titre\Large #2}\par\vspace{3pt}%
  {\color{encre}\hrule height 2pt}\vspace{6pt}}

% Encadrés : coins carrés, sans ombre
\newtcolorbox{regle}{enhanced,before upper=\raggedright,sharp corners,boxrule=0pt,leftrule=4pt,colback=papier,colframe=rouge,left=10pt,right=8pt,top=6pt,bottom=6pt}
\newtcolorbox{cadre}{enhanced,sharp corners,boxrule=2pt,colback=white,colframe=encre,left=10pt,right=10pt,top=8pt,bottom=8pt}

\setlist[itemize]{leftmargin=1.3em,topsep=2pt,itemsep=2pt,parsep=0pt}
\setlist[enumerate]{leftmargin=1.6em,topsep=2pt,itemsep=3pt,parsep=0pt}
\definecolor{codefond}{HTML}{EEEEE9}
% Code en ligne : fond gris pâle, comme `code` dans les diapos.
\newcommand{\code}[1]{{\setlength{\fboxsep}{1.6pt}\colorbox{codefond}{\ttfamily\small #1}}}
% Bloc de code : fond gris pâle, cadre encre, filet rouge à gauche (bloc-code des diapos).
\newtcolorbox{codebloc}{enhanced,sharp corners,colback=codefond,colframe=encre,boxrule=1pt,left=10pt,right=8pt,top=4pt,bottom=4pt,before skip=4pt,after skip=6pt,fontupper=\ttfamily\small,overlay={\fill[rouge] (frame.south west) rectangle ([xshift=3.5pt]frame.north west);}}
\newcommand{\bloc}[1]{\begin{codebloc}#1\end{codebloc}}
% Partie du travail : numéro rouge, puis titre
\newcommand{\partie}[2]{\par\vspace{7pt}\Needspace{4\baselineskip}{\titre\large\color{rouge}#1}\hspace{0.6em}{\titre\large #2}\par\vspace{1pt}}

\begin{document}
\raggedright
% Pas de césure : les motifs anglais coupaient mal le français (capa-ble).
\hyphenpenalty=10000
\exhyphenpenalty=10000
\setlength{\parindent}{0pt}

% Texte : consignes H24 d'Adrien Cloutier, reprises telles quelles.
% Changements A26 seulement : banque de données (CES 2025), Positron, monPortail,
% date, points-virgules retirés, coquilles corrigées, règle sur l'IA (plan de cours).

% ---------------------------------------------------------------- en-tête
\noindent\begin{minipage}[b]{0.7\textwidth}
{\titre\small\color{gris}POL-2000 · MÉTHODOLOGIE QUANTITATIVE · AUTOMNE 2026}\par\vspace{4pt}
{\titre\fontsize{24}{28}\selectfont Travail de mi-session}\par\vspace{2pt}
{\titre\large Consignes}
\end{minipage}\hfill
\begin{minipage}[b]{0.25\textwidth}\raggedleft
\includegraphics[height=0.55in]{ulaval-logo.png}
\end{minipage}
\par\vspace{6pt}{\color{rouge}\hrule height 2pt}\vspace{8pt}

\noindent\begin{tabularx}{\textwidth}{@{}*{4}{>{\raggedright\arraybackslash}X}@{}}
{\titre\small\color{gris}PONDÉRATION} & {\titre\small\color{gris}MODE} & {\titre\small\color{gris}REMISE} & {\titre\small\color{rouge}À REMETTRE}\\
{\titre\Large\color{rouge}20 \%} & Individuel & Format PDF, sur monPortail & \textbf{Mercredi 28 octobre, 23h59}\\
\end{tabularx}

% ---------------------------------------------------------------- description
\etape{DESCRIPTION}{Description}

Le travail de mi-session est une version courte de votre travail final. Il se veut une première version simplifiée de votre travail, à l'image d'un devis de recherche, pour obtenir une rétroaction du professeur avant la remise finale.

À l'intérieur de ce travail, vous devrez :
\begin{enumerate}
\item Concevoir une question de recherche mettant en relation une variable dépendante et une variable indépendante.
\item Présenter une courte revue de la littérature avec des théories permettant de dériver des hypothèses testables empiriquement.
\item Formuler une ou des hypothèses (et leur hypothèse nulle) sur la relation entre les deux variables.
\item Présenter la base de données et les variables sélectionnées pour tester les hypothèses.
\item Justifier la pertinence et l'originalité de la question de recherche.
\item Ajouter la contribution souhaitée à la littérature scientifique.
\item Conclure par une justification personnelle du choix de la recherche.
\end{enumerate}

Ceci est un travail individuel et aucune collaboration n'est permise. \textbf{Tou.te.s les étudiant.e.s devront travailler avec la banque de données de l'Étude électorale canadienne 2025 (la CES).} Cette banque de données et son codebook sont déjà dans votre dossier \code{pol2000/data} depuis l'examen 1. Sinon, vous pouvez les télécharger dans la console de Positron :
\bloc{ces::download\_ces\_dataset("2025")\\ces::download\_pdf\_codebook("2025")}

Le travail est à remettre en \textbf{format PDF} sur monPortail au plus tard le mercredi 28 octobre 2026, à 23h59. Veuillez prendre note qu'il n'est pas nécessaire d'utiliser Positron pour ce premier travail. Par contre, vous pouvez déjà prendre de l'avance sur le travail final, et réaliser de premiers graphiques descriptifs de vos données. Plus votre travail sera complet et avancé, plus vous aurez de la rétroaction du professeur en vue du travail final.

La présentation et la citation des sources seront prises en compte dans votre note. Consultez et respectez le guide de présentation des travaux du Département de science politique de l'Université Laval.

% ---------------------------------------------------------------- contenu
\etape{CONTENU}{Contenu}

À partir de la banque de données fournie, vous devrez créer un travail qui inclut les éléments suivants :

\partie{1}{Une question de recherche précise}
À quelle question voudriez-vous répondre ? Cette partie est primordiale. Pensez-y bien. Il faut qu'il soit possible de répondre à cette question avec les données à votre disposition. La question n'a pas besoin d'être compliquée. Elle doit mettre en relation une variable dépendante et une variable explicative (indépendante). Par exemple, je pourrais formuler la question suivante : le fait d'avoir des enfants (variable explicative, X) influence-t-il l'orientation gauche-droite des parents (variable dépendante, Y) ?

\partie{2}{Une hypothèse bien formulée, dérivée de la littérature}
Énoncez clairement votre hypothèse. Quelle sera, selon vous, la réponse à votre question ? Pourquoi ? Énoncez également votre hypothèse nulle en conséquence de votre hypothèse alternative. Cette section exige une mobilisation de la recherche existante pertinente. Je vous demande de rédiger minimum deux paragraphes incluant quatre références scientifiques qui permettent de soutenir votre hypothèse. Ces références doivent se retrouver dans la bibliographie du travail. Dégagez une (ou des, mais une serait plus simple dans le cadre de ce travail) théorie de votre littérature. Présentez-la, et montrez comment, à partir de cette théorie, vous pouvez dériver votre hypothèse.

\partie{3}{Une présentation des données}
Quelles variables voudrez-vous utiliser ? \textbf{Votre variable dépendante doit prendre une des unités de mesure suivantes pour être acceptée dans ce travail : ratio, intervalle et ordinal.} Autrement dit, des variables dépendantes nominales (dichotomiques ou autres) ne seront pas acceptées. Il est de votre entière responsabilité de vous assurer que la banque de données inclut les variables nécessaires pour répondre à votre question de recherche. Vous devez donc déterminer ceci en consultant la banque de données (et le codebook). Présentez vos données, votre base de données, pour démontrer votre connaissance et votre compréhension de celles-ci.

\partie{4}{Pertinence, contribution et justification personnelle}
En quoi votre question est-elle originale ? Comment souhaitez-vous contribuer à la littérature ? Quelle.s limite.s de la littérature souhaitez-vous démontrer, et combler ? Enfin, plus personnellement, expliquez-nous pourquoi vous avez choisi ce sujet. Est-il important pour vous ? Comment percevez-vous votre rôle de chercheur/chercheuse ? Ces réflexions sont au cœur du travail en recherche. Nous devons constamment nous questionner sur notre rôle, sur l'impact de notre travail, sur nos valeurs, sur nos biais et nos désirs de contribution. Laissez-vous aller dans votre réflexion !

% ---------------------------------------------------------------- format
\etape{FORMAT}{Format}

Le format du travail : longueur maximale de 2 pages (la page couverture et la bibliographie ne sont pas comprises), taille de police 12 pt, Times New Roman à 1,5 interligne. Le travail devra respecter le guide de présentation des travaux du Département de science politique de l'Université Laval. Les normes de citations et de références devront respecter le format APA (auteur-date).

\begin{regle}
\textbf{Intelligence artificielle.} Permis, à condition de le déclarer : obtenir de l'aide pour écrire, comprendre ou déboguer du code R. Non permis : l'interprétation des résultats et la rédaction des travaux. Ce que vous affirmez, vous devez pouvoir le défendre. À la fin du travail, après la bibliographie, déclarez si vous avez utilisé un outil d'IA. Si oui, lequel, comment et pourquoi ?
\end{regle}

\end{document}
